\scanpage{64}
\mspage{8}
% archive sheet number: 63
\diagram{\begin{tikzcd}[ampersand replacement=\&,column sep=2pt,row sep=2.6ex,cells={font=\scriptsize}]
\&\underline{\mathrm{Fibsc}}_0\,S\arrow[dl,hook,"\alpha'"']\arrow[d]
\&\&\underline{\mathrm{Cofibsc}}_0\,S\arrow[d]\arrow[dr,hook,"\beta'"]\&\\
\underline{\mathrm{Fibsc}}\,S\arrow[d]\&
\underline{\mathrm{Fib}}_0\,S\arrow[dl,hook,"\alpha"']\arrow[ddr,dashed,"\underline{\unsure{\mathrm{Lis}|\mathrm{Liss}}}_0"']
\&{}\&\mathrm{Cofib}_0\,S\arrow[dr,hook,"\beta"]\arrow[ddl,dashed,"\underline{\mathrm{Prop}}_0"]
\&\underline{\mathrm{Cofibsc}}\,S\arrow[d]\\
\underline{\mathrm{Fib}}\,S\arrow[dr,hook]\&
\&\underline{\unsure{\mathrm{Parf}|\mathrm{Perf}}}\,S\arrow[dl,hook]\arrow[d,hook]\arrow[dr,hook]
\arrow[u,phantom,"\deleted{\unsure{\Phi}}"]
\&\&\mathrm{Cofib}\,S\arrow[dl,hook]\\
\&\underline{\mathrm{Liss}\,S}\arrow[dr,hook]\&
\underline{\mathrm{Fibhot}}\,S\arrow[d,hook]\arrow[uuul,dashed,bend left=6,"\unsure{\Phi_0}"']\arrow[uuur,dashed,bend right=6,"\unsure{\Psi_0}"]\&
\underline{\mathrm{Prop}}\,S\arrow[dl,hook]\&\\
\&\&\mathrm{Cat}_{/S}\arrow[uuull,dashed,bend left=50,"\Phi"]\arrow[uuurr,dashed,bend right=50,"\Psi"']\&\&
\end{tikzcd}}

Je suis quasiment sûr du

\uline{Théorème}. Tous les foncteurs précédents induisent
des équivalences entre les catégories localisées.
% slanted
\margin{C'est totalement \unsure{canulé|coulé|croulé}, cf mes commentaires pages 32 ff. Solution des perplexités pages 41 ff. ! Voir \unsure{surtout} le ``théorème-scholie'' récapitulatif p. 70-81.}

Comme il y a 16 foncteurs, ça a l'air
d'un exercice de patience. Le principe général
est qu'on arrive à définir des foncteurs en
sens inverse, \inserted{compatibles avec les $W_S$,} de telle façon que les composés
soient \deleted{\unsure{iso}} reliés aux foncteurs identiques par
des flèches qui sont dans $W_S$, de sorte
que \deleted{par} par passage aux localisés ces foncteurs
induisent des (équivalences) quasi-inverses des
foncteurs \unsure{déduits|déduit} de ceux du diagramme.
Tous ces foncteurs en sens opposé résultent,
essentiellement par composition, de trois foncteurs
bien familiers. Deux de ceux-ci sont
indiqués en pointillés sur le diagramme, comme $\Phi,\Psi$.
